\section{问题分析}

三个问题共享相同的二体轨道动力学与单次光学观测几何条件，但决策结构逐层递进：从单目标两脉冲机动，扩展到单星多目标巡游排序，再进一步耦合双星任务分配与新星初始轨道设计。本节先分析单次观测的几何结构，再分析同步带附近轨道机动的代价特点，最后据此给出三问的建模思路。

\subsection{观测几何分析}

观测条件可以用目标周围的空间区域来理解。距离约束要求观测卫星位于以目标为球心、半径分别为 10~km 和 50~km 的两个球面之间。光照角约束要求观测卫星位于以目标指向太阳的射线为轴、半顶角为 $45^\circ$ 的圆锥内。侦测角约束要求观测卫星位于以目标指向地心的射线为轴、半顶角为 $45^\circ$ 的圆锥内。三者同时满足时观测卫星才处在可观测位置，观测卫星 R、目标卫星 B 与太阳、地心之间的几何关系如图~\ref{fig:geometry}(a) 所示。两个圆锥的轴随太阳和目标的运动不断转动，可观测区域因而是时变的。距离下界相当于在目标周围挖去了一个半径 10~km 的小球，区域中间是空的，不是凸集。这类区域上的优化问题没法直接按凸优化求解，第 6 节构造凸子问题时会用一个平面把这个空洞替换掉。

两个圆锥有公共部分的充要条件是两条轴线的夹角不超过 $90^\circ$。记目标处太阳方向与地心方向的夹角为 $\beta$，则可观测的必要条件为
\begin{equation}
  \beta(t) = \angle\big(\bm r_\odot(t) - \bm r_b(t),\ -\bm r_b(t)\big) \le 90^\circ .
  \label{eq:beta}
\end{equation}
太阳距离远大于同步轨道半径，这个条件近似等价于目标位于地球的背阳一侧。同步轨道目标每天绕地球一周，$\beta$ 也以一天为周期变化，每天大约有一半时间满足这个条件。图~\ref{fig:geometry}(b) 给出了目标 B1 在前 48~h 内 $\beta$ 的变化，阴影部分为满足必要条件的时段。这个条件只与目标和太阳有关，可以在求解之前预先算出每颗目标的可观测时段。

\begin{figure}[!htb]
  \centering
  \begin{subfigure}[b]{0.41\textwidth}
    \centering
    \includegraphics[width=\linewidth]{fig1a_geometry.pdf}
    \caption{观测几何关系}
  \end{subfigure}\hfill
  \begin{subfigure}[b]{0.57\textwidth}
    \centering
    \includegraphics[width=\linewidth]{fig1b_beta_b1.pdf}
    \caption{B1 处 $\beta$ 在前 48~h 内的变化}
  \end{subfigure}
  \caption{观测几何与可能观测时段}
  \label{fig:geometry}
\end{figure}

轨道面的差异进一步缩小了观测机会。R1 的轨道倾角为 $4.38^\circ$，30 颗目标的倾角分布在 $0.12^\circ$ 到 $9.01^\circ$ 之间。两颗卫星的轨道面相差 $1^\circ$ 时，法向相对运动的振幅约为 736~km，远大于 50~km 的观测距离上限。而在同步轨道上把轨道面转动 $1^\circ$ 约需 53.7~m/s 的速度增量，用 300~m/s 的预算逐一对准 30 颗目标的轨道面是不经济的。因此在低燃料的巡游方案中，观测卫星应优先利用两条轨道交线附近的自然近距离交会来完成观测。轨道交线每天经过两次，再叠加 $\beta \le 90^\circ$ 的条件，对于轨道面差异明显的目标，通常每天存在约一次低代价的观测机会。这一规律用来生成候选观测时刻，从而减少需要搜索的时段；求解时速度增量的方向不受限制，实际观测时刻由精确的轨道优化确定。

\subsection{机动代价分析}

所有卫星都在同步带附近，彼此之间的相对运动以沿轨道方向的相位漂移为主。设观测卫星的半长轴比同步轨道大 $\Delta a$，它相对同步轨道的经度漂移率为 $\dot\lambda = -\tfrac32 n\,\Delta a/a$，而在近圆轨道上改变半长轴 $\Delta a$ 需要的切向速度增量为 $\Delta v = \tfrac{v}{2a}\Delta a$。两式合起来得到
\begin{equation}
  \Delta v = \frac{v}{3n}\,\big|\Delta\dot\lambda\big| \approx 2.84~\mathrm{m/s}\ \text{每}\ 1^\circ/\text{天} .
  \label{eq:drift}
\end{equation}
要在 $D$ 天内追上相位相差 $\Delta\lambda$ 的目标，需要先建立 $\Delta\lambda/D$ 的漂移率，到达后再把它消除，代价约为 $5.68\,\Delta\lambda/D$~m/s。可见机动代价主要由相位差和可用时间决定。时间越充裕，所需的漂移率越小，代价也越低。

\begin{figure}[!htb]
  \centering
  \includegraphics[width=0.8\textwidth]{fig2_targets.pdf}
  \caption{初始时刻 30 颗目标在同步带上的分布（径向放大显示）}
  \label{fig:targets}
\end{figure}

图~\ref{fig:targets} 画出了初始时刻 30 颗目标在相位和半长轴上的分布。目标大致对称地分布在 R1 的东西两侧，各自延伸约 $19^\circ$。它们的半长轴都非常接近同步轨道，彼此漂移得很慢，一个月下来相对位置几乎不变。所以可以把这 30 颗目标看成一条基本固定的队列，R1 要做的是沿着这条队列依次走过去。由此可以得到两点判断。第一，访问顺序应沿相位单向推进，来回折返会让漂移的路程和燃料都成倍增加。第二，相位接近、轨道面也接近的目标可以在同一次飞越中先后观测，比如 B19 和 B20 几乎挨在一起，观测完一颗顺路就能观测另一颗，省去一次单独的转移。


\subsection{建模思路}

基于上述分析，三个问题可以沿“单次观测单元、单星序列巡游、双星选轨与分工”这条主线逐层建模，总体框架如图~\ref{fig:framework} 所示。


问题一是基本单元。给定 R1 的初始状态和目标 B1，决策变量是两次脉冲的时刻、两个速度增量矢量和观测开始时刻，约束是连续 60~s 内每一时刻都满足观测条件，这是一个带半无限约束的非线性规划。48~h 内 B1 有两个候选观测时段，越晚的时段留给相位调整的时间越多，所需速度增量越小。由于可观测区域非凸，求解时需要从多个初值出发搜索。约束要求 60~s 内每一时刻都满足，优化时每隔 1~s 检查一次，得到方案后再每隔 0.05~s 复查一遍，确认相邻检查时刻之间也没有越界。

\begin{figure}[!htb]
  \centering
  \includegraphics[width=0.78\textwidth]{fig3_framework.pdf}
  \caption{总体框架}
  \label{fig:framework}
\end{figure}

问题二把基本单元串成序列。决策同时包括离散部分和连续部分。离散部分是选哪些目标、按什么顺序、在哪一天的机会观测；连续部分是每次脉冲的时刻和矢量。前后两段转移共享中间的轨道状态，一段的机动会改变下一段的起点，所以逐段独立优化得到的方案并不是整体最优。对此，本文在离散层基于三维 Clohessy--Wiltshire（CW）能量积分两脉冲解析代价构建物理状态动态规划，并在全排列空间做确定性变邻域搜索以确定最优访问序列与交会日期；在连续层则将 30 天计划的全部脉冲矢量与施加时刻一并纳入序列凸优化做全计划联合精修，同时以传统束搜索和大邻域搜索作为基线对照。

问题三在问题二的外面增加了 R2 的轨道设计和目标分工。题面所说的 R2 提前发射，意味着它在 $t_0$ 时刻已经在轨，偏心率小于 0.001 约束的是 $t_0$ 时刻的初始轨道，这个初始轨道由我们设计。$t_0$ 之后的每次机动都计入 R2 的速度增量，机动后的轨道不再受偏心率约束。题目没有给出 R2 半长轴的具体范围，只说明它是高轨卫星。我们引入参数 $\delta a$，要求 R2 的初始半长轴与同步轨道半长轴之差不超过 $\delta a$，用来刻画“位于同步带内”这一要求。目标在相位上分布于 R1 的两侧，由两颗卫星各负责一侧并单向扫描，每颗卫星需要跨越的相位范围约减少一半。$\delta a$ 的取值和分工方案在第 7 节讨论。
